\par\medskip
\SourceLeaf{072}\pdfbookmark[2]{54.}{article-054}
\begin{center}54.\end{center}
Die in den vorangehenden Artikeln erläuterten Beziehungen der Bewegung zur Ekliptik gelten offenbar ebenso, wenn man die Ekliptik durch eine beliebige andere Ebene ersetzt, sofern die Lage der Bahnebene zu dieser Ebene bekannt ist. Die Ausdrücke Länge und Breite müssen dann entfallen. Damit stellt sich die Aufgabe, \textit{aus der bekannten Lage der Bahnebene und einer weiteren Ebene zur Ekliptik die Lage der Bahnebene zu dieser neuen Ebene abzuleiten.} Es seien $n\AscendingNode$, $\AscendingNode\AscendingNode'$ und $n\AscendingNode'$ Bögen der Großkreise, die von der Ebene der Ekliptik, der Bahnebene und der neuen Ebene auf der Himmelskugel erzeugt werden (Abb. 2). Um die Neigung des zweiten Kreises gegen den dritten und die Lage des aufsteigenden Knotens eindeutig angeben zu können, muss man auf dem dritten Kreis eine der beiden Richtungen als Entsprechung zu der Richtung wählen, die auf der Ekliptik der Reihenfolge der Tierkreiszeichen folgt. In unserer Abbildung sei dies die Richtung von $n$ nach $\AscendingNode'$. Ferner muss man von den beiden Halbkugeln, die der Kreis $n\AscendingNode'$ trennt, die eine als der nördlichen, die andere als der südlichen Halbkugel entsprechend ansehen. Die beiden Halbkugeln sind bereits dadurch unterschieden, dass stets diejenige als nördlich gilt, die rechts liegt, wenn man auf dem Kreis in der Reihenfolge der Tierkreiszeichen fortschreitet\footnote{Dabei ist die \textit{innere} Oberfläche der Kugel gemeint, die unsere Abbildung darstellt.}. In unserer Abbildung sind demnach $\AscendingNode$, $n$ und $\AscendingNode'$ die aufsteigenden Knoten des zweiten Kreises auf dem ersten, des dritten auf dem ersten und des zweiten auf dem dritten. Die Neigungen des zweiten gegen den ersten, des dritten gegen den ersten und des zweiten gegen den dritten sind der Reihe nach $180^0-n\AscendingNode\AscendingNode'$, $\AscendingNode n\AscendingNode'$ und $n\AscendingNode'\AscendingNode$. Unsere Aufgabe hängt somit von der Lösung eines sphärischen Dreiecks ab, dessen übrige Stücke aus einer Seite und den anliegenden Winkeln abzuleiten sind. Die gewöhnlichen Regeln, die in der sphärischen Trigonometrie für diesen Fall angegeben werden, lassen wir als hinreichend bekannt beiseite. Zweckmäßiger ist ein anderes Verfahren, das auf einigen Gleichungen beruht, die man in unseren Lehrbüchern der Trigonometrie vergeblich sucht. Die folgenden Gleichungen werden wir häufig verwenden. Dabei bezeichnen $a$, $b$ und $c$ die Seiten des sphärischen Dreiecks und $A$, $B$ und $C$ die jeweils gegenüberliegenden Winkel.\HistoricalNoteDouble{Handschriftliche Bemerkung:] Haben zwei Dreiecke die beiden Seiten $b$ und $c$ gemeinsam, so gilt
\[
\begin{aligned}
\mathrm{I.}\quad&\frac{\sin\frac{1}{2}(a'-a)}{\sin\frac{1}{2}(A'-A)}=\frac{\sin b\sin\frac{1}{2}(C'+C)}{\cos\frac{1}{2}(B'-B)}=\frac{\sin c\sin\frac{1}{2}(B'+B)}{\cos\frac{1}{2}(C'-C)}\\[.5em]
\mathrm{II.}\quad&-\frac{\sin\frac{1}{2}(a'+a)}{\sin\frac{1}{2}(A'-A)}=\frac{\sin b\cos\frac{1}{2}(C'+C)}{\sin\frac{1}{2}(B'-B)}=\frac{\sin c\cos\frac{1}{2}(B'+B)}{\sin\frac{1}{2}(C'-C)}\\[.5em]
\mathrm{III.}\quad&-\frac{\sin\frac{1}{2}(a'-a)}{\sin\frac{1}{2}(A'+A)}=\frac{\sin b\sin\frac{1}{2}(C'-C)}{\cos\frac{1}{2}(B'+B)}=\frac{\sin c\sin\frac{1}{2}(B'-B)}{\cos\frac{1}{2}(C'+C)}\\[.5em]
\mathrm{IV.}\quad&\frac{\sin\frac{1}{2}(a'+a)}{\sin\frac{1}{2}(A'+A)}=\frac{\sin b\cos\frac{1}{2}(C'-C)}{\sin\frac{1}{2}(B'+B)}=\frac{\sin c\cos\frac{1}{2}(B'-B)}{\sin\frac{1}{2}(C'+C)}\cdot
\end{aligned}
\]
}
\par\Needspace{11\baselineskip}
\SourceLeaf{073}
\[
\mathrm{I.}\quad\frac{\sin\frac{1}{2}(b-c)}{\sin\frac{1}{2}a{,}}=\frac{\sin\frac{1}{2}(B-C)}{\cos\frac{1}{2}A}
\]
\[
\mathrm{II.}\quad\frac{\sin\frac{1}{2}(b+c)}{\sin\frac{1}{2}a}=\frac{\cos\frac{1}{2}(B-C)}{\sin\frac{1}{2}A}
\]
\[
\mathrm{III.}\quad\frac{\cos\frac{1}{2}(b-c)}{\cos\frac{1}{2}a}=\frac{\sin\frac{1}{2}(B+C)}{\cos\frac{1}{2}A}
\]
\[
\mathrm{IV.}\quad\frac{\cos\frac{1}{2}(b+c)}{\cos\frac{1}{2}a}=\frac{\cos\frac{1}{2}(B+C)}{\sin\frac{1}{2}A}\cdot
\]
Den Beweis dieser Aussagen müssen wir hier der Kürze halber übergehen. Ihre Gültigkeit lässt sich jedoch für Dreiecke, deren Seiten und Winkel sämtlich höchstens $180^0$ betragen, von jedem leicht überprüfen. Fasst man das sphärische Dreieck in größtmöglicher Allgemeinheit auf, sodass weder den Seiten noch den Winkeln irgendwelche Grenzen gesetzt sind (was viele bedeutende Vorteile bietet, jedoch einiger Vorbemerkungen bedarf), so kann es Fälle geben, in denen in sämtlichen vorangehenden Gleichungen das Vorzeichen zu wechseln ist. Die ursprünglichen Vorzeichen ergeben sich aber offensichtlich wieder, sobald man einen der Winkel oder eine der Seiten um $360^0$ vergrößert oder verkleinert. Man kann die angegebenen Vorzeichen daher stets unbedenklich beibehalten, gleichviel ob die übrigen Stücke aus einer Seite und den anliegenden Winkeln oder aus einem Winkel und den anliegenden Seiten zu bestimmen sind. Unsere Formeln liefern nämlich immer entweder die gesuchten Werte selbst oder Werte, die um $360^0$ von den wahren abweichen und ihnen somit gleichwertig sind. Eine ausführlichere Erläuterung dieses Gegenstands behalten wir uns für eine andere Gelegenheit vor. Vorerst lässt sich jedoch durch strenge Induktion, das heißt durch vollständige Aufzählung aller Fälle, leicht nachweisen, dass die Regeln, die wir sowohl zur Lösung unserer Aufgabe als auch bei anderen Gelegenheiten auf diese Formeln stützen werden, in allen Fällen allgemein gelten.
